\documentclass[10pt,a4paper]{article}
\usepackage{amsmath,amssymb,amsthm,graphicx}
\usepackage{xfrac}
\newenvironment{enumeratei}
{\bgroup\def\theenumi{\roman{enumi}}\def\theenumii{\arabic{enumii}}\begin{enumerate}}
{\end{enumerate}\egroup}

\hyphenation{addi-tion-al belong belongs denote exchan-ge exchan-ged exchan-ges}

\newcommand{\psfrac}[2]{\sfrac{(#1)}{#2}}

\newcommand\mto{\mathchoice{\longmapsto}{\mapsto}{\mapsto}{\mapsto}}
\let\oldbigsqcup\bigsqcup
\renewcommand{\bigsqcup}{\mathop{\textstyle\oldbigsqcup}\displaylimits}

\let\epsilon\varepsilon

\DeclareMathOperator{\Ess}{Ess}
\DeclareMathOperator{\Leb}{Leb}

\usepackage{tikz}
\usetikzlibrary{plotmarks}

\theoremstyle{plain}
\newtheorem{theorem}{Theorem}[section]
\newtheorem{lemma}[theorem]{Lemma}
\newtheorem{proposition}[theorem]{Proposition}
\newtheorem{corollary}[theorem]{Corollary}
\newtheorem{claim}{Claim}

\theoremstyle{definition}
\newtheorem{remark}[theorem]{Remark}
\newtheorem{example}[theorem]{Example}

\newcommand{\A}{\mathcal{A}}
\newcommand{\B}{\mathcal{B}}
\newcommand{\I}{\mathcal{I}}

\newcommand{\Z}{\mathbb{Z}}
\newcommand{\N}{\mathbb{N}}
\newcommand{\R}{\mathbb{R}}
\title{On a symmetric subinterval the conjugate of the transformation is the reflection of that subinterval}
\author{Przemyslaw Berk, Frank Trujillo}
\date{doi:10.5802/jep.302; CC BY4.0}
\begin{document}
\maketitle
\noindent\textit{Editorial note.} The counted statement is the author's labelled lemma environment at source0.tex lines 596-602; its proof environment opens at line 604, the line after the blank separator, so the pairing is by adjacency and the retained proof range 604-613 is the whole environment.\par\medskip
\noindent\textit{Editorial note.} The numbered identity eq: centersandinvolution (source lines 519-521) is retained as context because the counted proof applies it at source line 607; it is the only outside label of the retained proof.\par\medskip
\noindent\textit{Editorial note.} This source is new to the material-authoring work and its preamble is built from the author's own macro and theorem block (source lines 8-45) rather than from a prepared adapter.\par\medskip
\noindent\textit{Editorial note.} This source has no label-number dump and the retained counter values are left empty, so the printed numbering of the card is the wrapper's default; the identity of the counted item is its author label.\par\medskip
\noindent\textit{Editorial scope.} One original result of the pinned publisher article, counted once: the conjugate of the transformation is the reflection of any symmetric subinterval (author label lem: local\_reflection).\par\medskip
\noindent\textit{Allowed standard prerequisites (editorial).}\begin{itemize}
\item Symmetric interval exchange transformation. An IET is symmetric when the reflection of the whole interval conjugates it to its inverse off the discontinuities.. It is the hypothesis.
\item Symmetric subinterval. A subinterval is alpha-symmetric when it is centred at the middle point c\_alpha of the exchanged interval I\_alpha.. It is the class of intervals on which the computation is performed.
\item Reflection of the centres. The conjugate of the transformation by the reflection of the whole interval fixes the middle point of each continuity interval.. It is the identity that converts the shifted centre back into the centre.
\end{itemize}\medskip
\begin{equation}\label{eq: centersandinvolution}
\mathcal{I}_I\circ T(c_{\alpha})=c_{\alpha},
\end{equation}

\begin{lemma}\label{lem: local_reflection}
Let $T : I \to I$ be a symmetric IET and $J \subseteq I$ be an $\alpha$-symmetric interval, for some $\alpha \in \A$. Then $\mathcal{I}_I\circ T(\mathring{J}) = \mathring{J}$. Moreover,
\[
\I_J = \mathcal{I}_I\circ T|_{\mathring{J}},
\]
that is, $\mathcal{I}_I\circ T|_{\mathring{J}}$ is the symmetric reflection on $J$.
\end{lemma}

\begin{proof}
Let $\alpha \in \A$ and $\Delta > 0$ such that $J = [c_\alpha - \Delta, c_\alpha + \Delta) \subseteq I_\alpha$. Fix $x\in J$ and let $-\Delta < \delta < \Delta$ be such that $x = c_{\alpha}+\delta$. Notice that $\I_J(x) = c_\alpha - \delta$.

Since $T|_{I_{\alpha}}$ is a translation, by~\eqref{eq: centersandinvolution}, we~have
\[
\I_I\circ T(x)=\I_I\circ T(c_{\alpha}+\delta)=\I_I\left(T(c_{\alpha}) +\delta\right)=
\I_I\circ T(c_{\alpha})-\delta=c_{\alpha}-\delta\,\in\, J,
\]
which finishes the proof.
\end{proof}

\end{document}
